% fees.tex — Fee structure and replication compatibility
% Kontrol: prove_cfmm_replication_feePreservation

\section{Fee Structure}\label{sec:fees}

\subsection{Fee Model}

The total trading fee is:
\begin{equation}\label{eq:total-fee}
  \fee(t) = \feebase + \feefund(t)
\end{equation}
where:
\begin{itemize}
  \item $\feebase$ is the base fee, fixed at deployment;
  \item $\feefund(t)$ is the funding fee component, updated at each liquidity event
        (see \cref{sec:funding-rate}).
\end{itemize}

\subsection{Fee Application}

Fees are applied per-swap on the input token. For a swap of $\Delta x$ input:
\begin{equation}\label{eq:fee-application}
  \Delta x_{\text{effective}} = (1 - \fee) \cdot \Delta x
\end{equation}
The fee amount $\fee \cdot \Delta x$ accrues to fee accumulators $\Fx$, $\Fy$
separately from the reserve invariant.

\subsection{Replication Preservation}

\begin{theorem}[Zero Replication Error]\label{thm:zero-error}
Under the replication-preserving fee model, the replication error is $\varepsilon = 0$.
\end{theorem}

\begin{proof}
The key property is that fees accrue \emph{separately} from the trading function invariant.

Let $(x, y)$ be the current reserves satisfying $\psi(x, y) = 0$.

\textbf{Step 1: Fee extraction.}
A swap submits $\Delta x$ input. The fee $f = \fee \cdot \Delta x$ is extracted
and added to the fee accumulator: $\Fx \gets \Fx + f$.
The effective input is $\Delta x_{\text{eff}} = \Delta x - f$.

\textbf{Step 2: Invariant-preserving swap.}
The output $\Delta y$ is determined by:
\begin{equation}\label{eq:invariant-swap}
  \psi(x + \Delta x_{\text{eff}},\; y - \Delta y) = 0
\end{equation}
Expanding with $\psi(x,y) = y + \ln(x) + 1 - x$:
\[
  (y - \Delta y) + \ln(x + \Delta x_{\text{eff}}) + 1 - (x + \Delta x_{\text{eff}}) = 0
\]

\textbf{Step 3: Replication check.}
Since the invariant $\psi = 0$ is maintained exactly before and after the swap
(fees do not enter the invariant equation), the reserve curve
$(x(p), y(p))$ continues to replicate $\VLP(p) = \ln(1+p)$ exactly.

The LP value at any price $p$ is:
\[
  \text{LP\_value}(p) = p \cdot x(p) + y(p) = p \cdot \frac{1}{1+p} + \ln(1+p) - \frac{p}{1+p} = \ln(1+p) = \VLP(p)
\]

The error $|\text{LP\_value}(p) - \VLP(p)| = 0$ for all $p$. \qedhere
\end{proof}

\begin{remark}[Compatibility with Bartoletti Gain Formula]
Since fees accrue separately, the Bartoletti gain formula \cite{bartoletti2024formalizing}
applies directly:
\[
  \text{gain}(\Gamma, t) = S_t \cdot r(\Gamma, t) - c_t
\]
where $r(\Gamma, t)$ is the swap rate computed from the fee-free invariant,
and $c_t$ includes the fee revenue as additional LP income. The fee-adjusted gain
is $\text{gain}_\fee = \text{gain} + \text{cumulative\_fees}$, which is always
at least as large as the fee-free gain.
\end{remark}

\subsection{Fee Bounds}

\begin{equation}\label{eq:fee-bounds}
  0 \leq \fee(t) \leq \fee_{\max}
\end{equation}
where $\fee_{\max}$ is a deployment-time constant. This ensures:
\begin{itemize}
  \item $\fee(t) \geq 0$: no negative fees (LP never pays trader);
  \item $\fee(t) \leq \fee_{\max}$: bounded fee prevents griefing.
\end{itemize}

The base fee satisfies $\feebase \leq \fee_{\max}$, and the funding component
satisfies $|\feefund(t)| \leq \fee_{\max} - \feebase$ (see CFMM-27).
